\begin{tikzpicture}[>=Stealth, thick]
    % 地面及阴影
    \draw (-3, 0) -- (4, 0);
    \foreach \x in {-2.8, -2.4, -2.0, -1.6, -1.2, -0.8, -0.4, 0, 0.4, 0.8, 1.2, 1.6, 2.0, 2.4, 2.8, 3.2, 3.6} {
        \draw (\x, 0) -- (\x-0.2, -0.2);
    }
    
    % 定义角度和圆心距
    \def\theta{20} % 倾角
    \def\d{4.5}    % 圆心距
    \def\r{1}      % 左侧小圆半径
    \def\R{1.2}    % 右侧大圆半径
    
    % 计算两圆心坐标
    \coordinate (O1) at (0, { \r / cos(\theta) });
    \coordinate (O2) at (\theta:\d);
    \coordinate (O2) at ($(O1) + (\theta:\d)$); % 实际 O2 位置
    
    % 绘制支撑杆 (简化为折线)
    \draw (O2) -- (2.5, 0);
    
    % 绘制两个圆
    \draw (O1) circle (\r);
    \draw (O2) circle (\R);
    \fill (O1) circle (1.5pt);
    \fill (O2) circle (1.5pt);
    
    % 绘制相切的皮带 (外公切线)
    % 倾角调整切点计算
    \pgfmathsetmacro{\alpha}{asin((\R-\r)/\d)}
    \coordinate (T1top) at ($(O1) + (90+\theta+\alpha:\r)$);
    \coordinate (T2top) at ($(O2) + (90+\theta+\alpha:\R)$);
    \coordinate (T1bot) at ($(O1) + (-90+\theta-\alpha:\r)$);
    \coordinate (T2bot) at ($(O2) + (-90+\theta-\alpha:\R)$);
    
    \draw (T1top) -- (T2top);
    \draw (T1bot) -- (T2bot);
    
    % 绘制底部的斜面延长线以标记角度
    \draw (-2, 0) -- (O1); % 简化示意，实际是相切到底部水平线
    % 更准确的底边斜线
    \coordinate (BottomIntersect) at (-1.5, 0);
    \draw (BottomIntersect) -- ($(BottomIntersect) + (\theta:3.5)$);
    
    % 角度 \theta
    \draw (-0.5, 0) arc (0:\theta:1);
    \node at (-0.2, 0.2) {$\theta$};
    
    % 标签
    \node[left] at ($(O1) + (-1.1, 0)$) {$m, r$};
    \node[right] at ($(O2) + (1.3, 0)$) {$M, r$}; % 原图文字似为 M,r
\end{tikzpicture}